\section{文化与思想：经典、文字与多重权威}

汉代的“经学”不是朝廷单向灌输的一套固定教本。五经博士、太学、察举、师承和地方评价，使经典解释成为进入官署、建立声望和争论政务的渠道。白虎观、石经和官方礼仪能显示朝廷试图公开规定何种文字和解释可被援引，不能证明所有地方按同一课程学习，或所有士人有同一政治立场。

熹平石经的残石、碑刻、墓葬和传世经传，分别保存文本、赞助、仪式和精英生活的不同层面。它们相互照亮而不相互代写：一块残石不能说明全国读者，一部后出的传记也不能替代某次讲学的现场记录。党锢尤其提醒我们，知识网络既可为朝廷提供评价官员的资源，也可能被朝廷视为会形成独立动员的关系。

宗教与方术同样不能压成“儒家之外”。楚王英案中的黄老、佛教供养文字是早期可见线索，不是全国已有统一佛教制度的证据；太平道等宗教语言在黄巾时可参与组织和解释危机，也不能把不同地区的参与者还原为单一信仰共同体。墓葬、符箓、早期佛教材料和正史记载的保存目的不同，足以显示汉人拥有多重求福、疗病、结社与解释秩序的途径，难以给出完整信众数目或私人动机。

写作本身是帝国的一项能力。诏令、簿籍、石经、碑铭和史书让命令、评价和记忆跨过距离；谁控制书写的格式，谁并不必然控制所有读者和执行者。班超、甘英所带回的远方信息，燕然山铭所作的功绩宣示，王莽和曹魏的受禅文辞，都应作为有立场的文字行动来读。它们不因立场而失去史料价值，却也不能被当作透明的事实录音。
